\subsection{محدودیت‌ها و تحلیل شکست}
تحلیل عوامل افت کارایی و محدودیت‌های الگوریتم حاکی از نکات زیر است:
\begin{itemize}
    \item \textbf{افت ناشی از صفر کردن سخت سطرها بدون بازآموزی (\lr{Hard Zero Pruning}):} صفر کردن وزن‌های سطرها بدون آموزش جبرانی موجب حذف برخی سیگنال‌های پروجکشن سرها شده و در مدل‌های بزرگ‌تری که حساسیت بالا به هماهنگی سرها دارند، دقت اولیه را کاهش می‌دهد.
    \item \textbf{انحراف در فشرده‌سازی لایه‌های دوطرفه:} در مدل‌های دوطرفه (\lr{Bidirectional Encoder}) مانند \lr{RoBERTa}، مفهوم چاه‌های توجه با مدل‌های زبانی خودبازگشت متفاوت است؛ چرا که توکن‌های پایانی نیز به صورت یکنواخت به توکن‌های پیشین متصل می‌شوند و اتکای صِرف به چاه‌های اولیه موجب انحراف میدان توجه می‌گردد.
    \item \textbf{وابستگی دقت به آستانه خطای فروبنیوس:} در صورتی که آستانه خطای مجاز بیش از اندازه آزاد در نظر گرفته شود (مانند $0.85$)، هرس شدید رخ می‌دهد و بازیابی دقت نیازمند استفاده از مکانیزم بازگشت به وزن‌های دست‌نخورده لایه خواهد بود.
\end{itemize}
